% ============================================================================
% Capítulo 1 — Caracterização da organização
% ============================================================================
\section{CARACTERIZAÇÃO DA ORGANIZAÇÃO}

\subsection{Segmento de atuação e perfil institucional}

A Cooperativa Agropecuária Vale Verde atua no segmento de comercialização de
produtos da agricultura familiar. Fundada em 2011 por um grupo de dezessete
famílias produtoras, a organização surgiu para viabilizar a venda conjunta da
produção e reduzir a dependência de intermediários. Atualmente, reúne 52
produtores associados distribuídos em quatro comunidades -- tradicional,
quilombola, assentamento e indígena -- em municípios da região metropolitana de
Campinas, no interior paulista.

O portfólio é organizado em três linhas: hortifrúti (tomate, alface, cenoura,
abobrinha e temperos verdes), grãos (feijão e milho) e produtos beneficiados
(polpas de fruta e farinha de mandioca). A produção é pulverizada em pequenas
propriedades, com média de 3,4 hectares por família, e a colheita ocorre em ciclos
curtos, o que exige coordenação diária entre oferta, demanda e logística.

A receita da cooperativa provém de três canais: compras públicas (PNAE e PAA),
fornecimento para restaurantes universitários e venda para redes varejistas
regionais. Os dois primeiros canais exigem comprovação documental rigorosa da
entrega e rastreabilidade do lote; o terceiro concentra a maior parte do
faturamento, porém impõe prazos de pagamento mais longos e maior pressão por
preço.

\subsection{Estrutura organizacional}

A governança segue o modelo cooperativista clássico, com assembleia geral como
instância máxima de deliberação. A diretoria eleita é composta por presidente,
vice-presidente e tesoureiro, apoiados por um conselho fiscal. A operação é
conduzida por uma equipe administrativa de quatro pessoas, um técnico agrícola
responsável pela assistência e pelo acompanhamento da qualidade, e três
motoristas que realizam as rotas de coleta e entrega em dois veículos próprios.

A Figura~\ref{fig:org} representa essa estrutura de forma simplificada. Observa-se
que a mesma equipe administrativa acumula funções de cadastro, atendimento,
faturamento e prestação de contas, o que explica parte dos gargalos
identificados na seção seguinte.

\begin{figure}[H]
  \centering
  \caption{Organograma simplificado da Cooperativa Vale Verde}
  \label{fig:org}
  \begin{tikzpicture}[
    node distance=8mm and 10mm,
    caixa/.style={draw, rounded corners=2pt, align=center, font=\fontsize{9}{11}\selectfont,
      minimum height=9mm, minimum width=30mm, inner sep=3pt},
    seta/.style={-latex, thin}
  ]
    \node[caixa, minimum width=44mm] (ag) {Assembleia geral};
    \node[caixa, below=of ag] (dir) {Diretoria\\ (presidente, vice, tesoureiro)};
    \node[caixa, below=of dir] (adm) {Equipe administrativa\\ (4 pessoas)};
    \node[caixa, below=of adm, text width=26mm] (prod) {Produtores\\ associados (52)};
    \node[caixa, right=of adm, text width=28mm] (tec) {Técnico agrícola\\ (1)};
    \node[caixa, left=of adm, text width=28mm] (log) {Logística\\ (3 motoristas)};
    \node[caixa, below=of adm, text width=30mm, yshift=-8mm] (fisc) {Conselho fiscal};
    \draw[seta] (ag) -- (dir);
    \draw[seta] (dir) -- (adm);
    \draw[seta] (dir.south west) -- (fisc.north);
    \draw[seta] (adm) -- (prod);
    \draw[seta] (adm) -- (tec);
    \draw[seta] (adm) -- (log);
    \draw[seta] (prod.east) to[bend left=18] (log.south);
  \end{tikzpicture}
  \fonte{elaborado pelo autor (2026)}
\end{figure}

\subsection{Principais processos}

O ciclo operacional da cooperativa pode ser descrito em oito processos de
negócio, apresentados na Tabela~\ref{tab:processos} com os problemas observados e
as oportunidades de melhoria por meio de tecnologia.

\begin{table}[H]
  \centering
  \caption{Processos de negócio, problemas e oportunidades}
  \label{tab:processos}
  \fontsize{10}{12}\selectfont
  \begin{tabular}{L{3.2cm}L{5.4cm}L{5.6cm}}
    \toprule
    \textbf{Processo} & \textbf{Problema atual} & \textbf{Oportunidade} \\
    \midrule
    Cadastro de produtor & Fichas em papel, dados incompletos e sem
      histórico de aprovação & Cadastro digital com aprovação e trilha de
      auditoria \\
    Planejamento da colheita & Comunicação por telefone e mensagens; oferta
      desconhecida & Registro de lote com quantidade e data de colheita \\
    Recebimento de pedidos & Pedidos em canais distintos, sem padrão de itens e
      preços & Pedido digital com itens, lote e endereço de entrega \\
    Separação e conferência & Conferência manual, sujeita a divergências &
      Itens vinculados ao lote e ao pedido, com total calculado \\
    Programação de rotas & Planilha sem vínculo com pedidos e veículos &
      Rota com data, motorista, veículo e entregas geradas \\
    Entrega e comprovação & Foto enviada por aplicativo de mensagens, sem
      vínculo com o pedido & Evidência com foto, GPS e observação no próprio
      registro \\
    Faturamento e repasse & Emissão manual, com atrasos e erros de valor &
      Fatura gerada a partir do pedido, com numeração sequencial \\
    Prestação de contas & Consolidação manual de planilhas & Relatórios e
      indicadores consolidados \\
    \bottomrule
  \end{tabular}

  \vspace{2pt}
  {\fontsize{10}{12}\selectfont\itshape Fonte: elaborado pelo autor (2026)}
\end{table}

\subsection{Problemas identificados}

A análise dos processos revelou quatro problemas estruturais. O primeiro é a
\textbf{descontinuidade da informação}: cada etapa usa um suporte diferente
(caderno, planilha, mensagem, nota fiscal), sem chave comum que permita
reconstruir o caminho entre a colheita e o consumo. O segundo é a
\textbf{baixa confiabilidade do estoque}: sem registro tempestivo de entradas e
saídas, a quantidade disponível de um lote diverge do que existe fisicamente,
provocando promessas de venda que não podem ser cumpridas.

O terceiro problema é a \textbf{fragilidade da comprovação de entrega}. A
fotografia enviada por aplicativo de mensagens não contém identificação do
pedido, coordenadas ou assinatura do recebedor, o que dificulta a resolução de
contestações e a prestação de contas em programas públicos. O quarto é a
\textbf{baixa visibilidade gerencial}: a diretoria não dispõe de indicadores
consolidados de vendas por produtor, perdas por lote ou desempenho logístico, e
decide com base em percepção e não em dados.

\subsection{Análise SWOT}

A Tabela~\ref{tab:swot} sintetiza o diagnóstico estratégico da organização,
considerando fatores internos e externos que condicionam a adoção da solução
tecnológica.

\begin{table}[H]
  \centering
  \caption{Matriz SWOT da Cooperativa Vale Verde}
  \label{tab:swot}
  \fontsize{10}{12}\selectfont
  \begin{tabular}{L{7.2cm}L{7.2cm}}
    \toprule
    \textbf{Forças} & \textbf{Fraquezas} \\
    \midrule
    Diversidade produtiva e vínculo com comunidades tradicionais; acesso a
    compras públicas; reputação regional; associativismo consolidado &
    Controles manuais; ausência de rastreabilidade; dependência de poucas
    pessoas na administração; baixa maturidade digital da base \\
    \midrule
    \textbf{Oportunidades} & \textbf{Ameaças} \\
    \midrule
    Digitalização de compras públicas; demanda por produtos com origem
    comprovada; programas de crédito para agricultura familiar; expansão do
    canal varejista & Concorrência de distribuidores com logística estruturada;
    exigências sanitárias crescentes; variação de preços; exclusão digital de
    parte dos produtores \\
    \bottomrule
  \end{tabular}

  \vspace{2pt}
  {\fontsize{10}{12}\selectfont\itshape Fonte: elaborado pelo autor (2026)}
\end{table}

